% !PRINTED 250
\weberpage{250}
From this there follows, as above, the conclusion:

1. \gesperrt{For every $n$ the quantities}
\[
P^{24}_{c,\eth,a}
\]
\gesperrt{are the roots of a transformation equation.}

If $n$ is not divisible by 3, one may assume the $c, c_2, c_1$ to be
divisible by 3. Then
\[
\lambda \equiv n, \quad \alpha \equiv 0, \quad \delta \equiv 0 \pmod{3}.
\]
Thus, as follows from §\,38, (15), $r$ is an \gesperrt{eighth root of unity};
if we therefore also observe
\[
\gamma_2(\omega+1) = e^{-\frac{2\pi i}{3}}\gamma_2(\omega), \quad
\gamma_2\left(-\frac{1}{\omega}\right) = \gamma_2(\omega), \quad \text{[§\,54, (14)]}
\]
we have the theorem:

2. \gesperrt{If $n$ is not divisible by 3 and $c$ is divisible by 3,
then the quantities}
\[
P^{8}_{c,\eth,a}\,\gamma_2(\omega)^{n-1}
\]
\gesperrt{are roots of a transformation equation.}

To make the application to the simplest case $n=2$, we set

\weq{3}{x = 16\,\frac{\eta(2\omega)^8}{\eta(\omega)^8}, \quad
x_0 = \frac{\eta\left(\dfrac{\omega}{2}\right)^8}{\eta(\omega)^8}, \quad
x_1 = \frac{\eta\left(\dfrac{\omega+3}{2}\right)^8}{\eta(\omega)^8},}

and then $x, x_0, x_1$ are nothing other than our functions
\[
f_2(\omega)^8, \quad f_1(\omega)^8, \quad -f(\omega)^8. \qquad \text{[§\,34, (9)]}
\]
These are, as we have already seen earlier [§\,54, (8)], the
roots of the cubic equation
\weq{4}{x^3 - x\,\gamma_2(\omega) + 16 = 0.}
The circumstance that $f_2^8, f_1^8, -f^8$ are themselves roots of a transformation equation for the transformation degree 2 explains the
phenomenon that on the adjunction of these quantities, and hence also on
the adjunction of $x^2$, the transformation equations belonging to an even $n$ become reducible.

If $n$ is an odd number, one takes $c$ to be \gesperrt{divisible by 4}; then
it follows that [§\,69, (9), (11), (12)]
\[
\lambda \equiv n \pmod{4},
\]
\[
\alpha \equiv 0, \quad \delta \equiv 0, \quad \beta \equiv a\,\eth_2, \quad \gamma \equiv -\eth\,a_2 \pmod{4},
\]
and hence there results
\[
r^6 = (-1)^{\frac{\beta-1}{2}} = (-1)^{\frac{n-1}{2}}(-1)^{\frac{a-a_2}{2}}
\]
%CONT>>
